\documentclass[12pt,a4paper]{article}
\usepackage[utf8]{inputenc}
\usepackage[english]{babel}



\title{Software Requirements Specification: Engine Control Demand Characterizer}
\date{\today}
\author{Aaron Willcock}


\begin{document}
\maketitle



\section{Introduction}
\subsection{Purpose}

This software aims to:
\begin{enumerate}
    \item Define a simplified engine model (e.g., accel, speed parameters),
    \item define a set of Adaptive Variable-Rate (AVR) tasks as proposed by
    Biondi et al. \cite{biondi} (e.g., transition speeds, execution times),
    \item combine the engine model and AVR task to form a real-time AVR-based
    engine control model, hereafter AEM (AVR Engine Model),
    \item calculate the maximum demand (see definition CITE)
    a given AEM can generate, and
    \item support result synthesis and visualization (e.g. defining experiment
    classes and parameters, exporting test data into portable files, enabling 
    comparing algorithm runtimes, etc.)
\end{enumerate}

The intended use is demand characterization for generated AEMs.
In furtherance of this intended use, the software also aims to
enumerate multiple calculation options (e.g., different methods).

\subsection{Scope}

The scope of this software includes three primary areas:
\begin{enumerate}
    \item model definition and generation,
    \item algorithm implementation,
    \item parallelization support, and
    \item synthesis and data visualization.
\end{enumerate}

For model definition and generation, scope includes:
\begin{enumerate}
    \item enumeration of existing models and
    \item randomized generation based on existing models.
\end{enumerate}
Excluded items include:
\begin{enumerate}
    \item creation of new engine models and
    \item parameter extension beyond existing models.
\end{enumerate}

For algorithm implementation, scope includes:
\begin{enumerate}
    \item implementation of existing algorithms and
    \item parameterization of algorithm variations where feasible
    (e.g., parameterization of constants in existing research).
\end{enumerate}

Excluded items include:
\begin{enumerate}
    \item design and implementation of new algorithms and
    \item approximation of algorithms no already specified.
\end{enumerate}

For synthesis and data visualization, scope includes:
\begin{enumerate}
    \item definition and implementation of experiment templates,
    \item parameterization of experiment variables,
    \item design and implementation of experiment test harness(es) for exporting
    key performance indicators (e.g., demand calculated, runtime, RAM use, etc.)
    , and
    \item definition and use of common key performance data export format (e.g.,
    CSV enumeration of key performance indicators).
\end{enumerate}
\subsection{Definitions, Acronyms, Abbreviations}

\begin{itemize}
    \item AVR - Adaptive Variable Rate - The Adaptive Variable Rate task model by
    Biondi et al. CITE
    \item HPC - High Performance Computing - High performance computing, usually a
    grid, cluster, or collection of compute devices
    \item ICE - Internal Combustion Engine - Heat engine using fuel combustion 
    to cause high-temperature, high-pressure gas expansion for performing work
    \item SI ICE - Spark Ignition ICE - Spark ignition ICE
\end{itemize}

\subsection{Acknowledgements}
This work templated from IEEE 830-1984 - IEEE Guide for Software Requirements
Specifications \cite{ieee_ieee_1984}.

\subsection{References}
\bibliographystyle{abbrv}
\bibliography{ecdc}

\subsection{Overview}
This software should accept as input:
\begin{enumerate}
    \item an ICE model,
    \item an AVR task model,
    \item an algorithm selection,
    \item a demand window.
\end{enumerate}
For the input above, this software should provide:
\begin{enumerate}
    \item the demand bound and
    \item the calculation runtime.
\end{enumerate}

\section{General Description}
\subsection{Product Perspective}

This product is aims to:
\begin{enumerate}
    \item implement existing demand characterizations and
    \item serve as a comparison point for future demand characterizers.
\end{enumerate}

\subsection{Product Functions}

The primary functions of this software include:
\begin{enumerate}
    \item enforcing invariants for ICE and AVR task models,
    \item viewing and selecting demand characterization algorithms,
    \item applying the above algorithms to the ICE and AVR models, and
    \item generating demand bounds.
\end{enumerate}

\subsection{User Classes and Characteristics}

Only one class of user exists with access to all features.

\subsection{Operating Environment}

This software is expected to run in a unix environment.

\subsection{Design, Implementation, and Environmental Constraints}

\subsubsection{Design Constraints}
The design should bias towards portability and abstraction.

\subsubsection{Implementation Constraints}
The implementation is restricted to:
\begin{enumerate}
    \item C++14 or newer, and
    \item a user-defined RAM limit.
\end{enumerate}
The C++14 requirement satisfies reviewer desires for a compiled implementation.
The RAM limit respects execution on standalone vs HPC systems.
\subsubsection{Environmental Constraints}
There are no known environmental constraints at this time.

\subsection{Assumptions and Dependencies}
\subsubsection{Technical}
This software assumes:
\begin{enumerate}
    \item a C++14 compiler is available for the end unix system and
    \item RAM limits comport with the executing system.
\end{enumerate}
\subsubsection{Operational}
This software assumes the following information is provided:
\begin{enumerate}
    \item ICE model parameters (min/max acceleration, min/max speed),
    \item AVR model parameters (mode count, transition speeds), and
    \item RAM limits
\end{enumerate}
\subsubsection{Business}
There are no known business assumptions or dependencies at this time.
\subsubsection{Environmental}
Uninterrupted system power for the duration of characterization is assumed.

\section{Specific Requirements}
System function requirements (FR) are enumerated below.
All FRs use the abbreviation USBAT, "users should be able to" and 
USBUT, "users should be unable to" brevity.
\subsection{FR 1: ICE Model CRUD}
USBAT:
\begin{enumerate}
    \item provide ICE max acceleration in UNITS,
    \item provide ICE min acceleration in UNITS,
    \item provide ICE symmetric acceleration in UNITS, and
    \item receive an ICE model data structure containing this information.
\end{enumerate}
USBUT:
\begin{enumerate}
    \item create invalid ICE models or accelerations.
\end{enumerate}

\subsection{FR 2: AVR Model CRUD}
USBAT:
\begin{enumerate}
    \item provide AVR task transition speeds in UNITS,
    \item provide AVR task worst-case execution times (WCETs) in UNITS, and
    \item receive an AVR task model data structure containing this information.
\end{enumerate}
USBUT:
\begin{enumerate}
    \item create invalid AVR models, transition speeds, or WCETs.
\end{enumerate}

\subsection{FR 3: AVRICE Model CRUD}
USBAT:
\begin{enumerate}
    \item provide an ICE model,
    \item provide AVR model, and
    \item receive an AVRICE model data structure containing this information.
\end{enumerate}
USBUT:
\begin{enumerate}
    \item create invalid AVRICE models from incompatible ICE and AVR models.
\end{enumerate}

For the following features, an AVRIACE DC algorithm is abbreviated as ADCA.
\subsection{FR 4: AVRICE Job CRUD}
USBAT:
\begin{enumerate}
    \item provide an AVRICE model,
    \item provide a ADCA, and
    \item receive an AVRICE Job containing this information.
\end{enumerate}
USBUT:
\begin{enumerate}
    \item create invalid AVRICE Jobs from incompatible AVRICE and ADCA models.
\end{enumerate}

\subsection{FR 5: ADCA implementation of 2017 Refinement of Workload algorithm (ROWL'17)}
USBAT:
\begin{enumerate}
    \item select KAVR'17 as the ADCA as a component of an AVRICE Job.
\end{enumerate}
\subsection{FR 6: ADCA implementation of 2018 Knapsack AVR (KAVR'18) CITE}
USBAT:
\begin{enumerate}
    \item select KAVR'18 as the ADCA as a component of an AVRICE Job.
\end{enumerate}
\subsection{FR 7: ADCA implementation of 2024 FPTAS Exact (XACT'24) CITE}
USBAT:
\begin{enumerate}
    \item select XACT'24 as the ADCA as a component of an AVRICE Job.
\end{enumerate}
\subsection{FR 8: ADCA implementation of 2024 FPTAS APX'24}
USBAT:
\begin{enumerate}
    \item select APX'24 as the ADCA as a component of an AVRICE Job.
\end{enumerate}

\subsubsection{FR 9: AVRICE Job Executor}
USBAT:
\begin{enumerate}
    \item provide an ordered sequence of AVRICE Jobs for execution and
    \item receive a corresponding, ordered sequence demand bounds and runtimes.
\end{enumerate}


\subsection{External Interface Requirements}

\subsubsection{User Interfaces}
No GUI is provided for UI.
Users interface through user-provided source code.

\subsubsection{Hardware Interfaces}
No hardware interfaces are required.

\subsubsection{Software Interfaces}
The library includes the following interfaces:

\begin{enumerate}
    \item ICE Model CRUD API,
    \item AVR Model CRUD API,
    \item AVRICE Model CRUD API,
    \item AVRICE Job CRUD API,
    \item AVRICE Executor Service API.
\end{enumerate}

\subsection{Non-Functional Requirements / Performance Requirements}
System non-functional requirements (FR) are enumerated below.
All FRs use the abbreviation SSB, "the system should be" and 
SSNB, "the system should not be", for brevity.

\subsubsection{NF1: Modularity}
SSB:
\begin{enumerate}
    \item distributed across single-purpose classes, and
    \item implemented with single-purpose, standalone services.
\end{enumerate}
\subsubsection{Availability}
SSB:
\begin{enumerate}
    \item able to sustain failed jobs without compromising other jobs,
    \item able to provide partial results for failed jobs,
    \item able to restrict jobs exceeding user-defined response times.
\end{enumerate}
\subsubsection{Concurrency}
SSB:
\begin{enumerate}
    \item able to launch AVRICE Jobs as individual processes,
    \item able to respect user-defined process limits.
\end{enumerate}
\subsubsection{Reliability}
SSB:
\begin{enumerate}
    \item able to gracefully handle invalid user inputs to any module and
    \item respectful of user-defined resource limits.
\end{enumerate}

\end{document}
